% !TeX program = lualatex
% ================================================================
% Chinese (Traditional) Mathematical Dictionary
%
% Generated by the server using the following settings:
%
%   language            Chinese (Traditional) (zht)
%   text direction      left to right
%   font                Noto Sans CJK TC
%   fallback font       Noto Serif
%   built from          latex/dictionary/mathematicalDictionary.tex
%   tier 3 vocabulary   green   RGB 0 128 0
%   English text        black   RGB 0 0 0
%   Chinese (Traditional) textpurple  RGB 112 48 160
%   layout macros       version 3
%
% Please don't edit any information in this comment block - the server
% compares it against the current settings and rebuilds this file when
% they differ, so changes here would be overwritten. However, feel free
% to make updates to the LaTeX below if you want to edit it.
% ================================================================
%
% ---- What is safe to edit in this file ----------------------------------
%
% Safe:     the Chinese (Traditional) word and the Chinese (Traditional) definition - the third
%           and fourth arguments of each entry. That is what this file is for:
%           if a translation reads wrongly to somebody who speaks the language,
%           correct it here and every sheet that uses the word follows.
%
% Not safe: the English word and the English definition - the first and second
%           arguments. The first is how an entry is matched to the shared
%           dictionary; the second is the wording the translation was made
%           from, and is how the server knows the translation is still current.
%           Changing an English definition here does not change the shared
%           dictionary - it only makes this entry look up to date when it is
%           not. Edit the English in the shared dictionary instead, and this
%           file will be brought back into step on the next run.
%
% Entries are added and removed by the server, following the shared dictionary.
% An entry the server cannot read is reported rather than guessed at, so if a
% run complains about this file, look for a missing brace.
% -------------------------------------------------------------------------

\documentclass[11pt]{article}
\usepackage[a4paper,margin=18mm]{geometry}
% Characters this language's font has not got are borrowed from another,
% rather than being silently dropped - see docs/translated-sheet-latex.md
\directlua{%
  if luaotfload and luaotfload.add_fallback then
    luaotfload.add_fallback("eallatinfallback",
      { "Noto Serif:mode=harf;" })
  else
    tex.error("this luaotfload is too old to register a fallback font")
  end
}
\usepackage[english]{babel}
\babelprovide[import]{chinese-traditional}
\babelfont[chinese-traditional]{rm}[Renderer=Harfbuzz, BoldFont={Noto Sans CJK TC}, ItalicFont={Noto Sans CJK TC}, BoldItalicFont={Noto Sans CJK TC}, RawFeature={fallback=eallatinfallback}]{Noto Sans CJK TC}
\babelfont[chinese-traditional]{sf}[Renderer=Harfbuzz, BoldFont={Noto Sans CJK TC}, ItalicFont={Noto Sans CJK TC}, BoldItalicFont={Noto Sans CJK TC}, RawFeature={fallback=eallatinfallback}]{Noto Sans CJK TC}
\newcommand{\ealtext}[1]{\foreignlanguage{chinese-traditional}{#1}}
\newcommand{\ealtextblock}[1]{{\raggedright\foreignlanguage{chinese-traditional}{#1}\par}}
% ================================================================
% EAL layout helpers
% ================================================================
\usepackage{xcolor}

\definecolor{ealkeycolour}{RGB}{0,128,0}
\definecolor{ealenglish}{RGB}{0,0,0}
\definecolor{ealtwocolour}{RGB}{112,48,160}

\newsavebox{\ealboxen}
\newsavebox{\ealboxtr}
\newlength{\ealglosswd}
\newlength{\ealglossraise}
\setlength{\ealglossraise}{1.15em}

% a tier 3 word, in English and in translation alike
\newcommand{\ealkey}[1]{\textcolor{ealkeycolour}{#1}}
\newcommand{\ealkeytr}[1]{\textcolor{ealkeycolour}{\ealtext{#1}}}

% One translated word above one English word. Built from kernel
% primitives, never a tabular, which would break wherever the sheet
% loads array, booktabs or tabularx. The box is as wide as the wider of
% the two, so a long translation pushes its neighbours aside instead of
% overlapping them, and it claims its full height so a table row or a
% TikZ node makes room for it.
\newcommand{\ealgloss}[2]{%
  \sbox{\ealboxen}{\ealkey{#1}}%
  \sbox{\ealboxtr}{\small\textcolor{ealtwocolour}{\ealtext{#2}}}%
  \setlength{\ealglosswd}{\wd\ealboxen}%
  \ifdim\wd\ealboxtr>\ealglosswd\setlength{\ealglosswd}{\wd\ealboxtr}\fi
  \leavevmode
  \makebox[\ealglosswd][c]{%
    \makebox[0pt][c]{\usebox{\ealboxen}}%
    \makebox[0pt][c]{\raisebox{\ealglossraise}{\usebox{\ealboxtr}}}%
  }%
}

% opens up line spacing around a block of glosses, so the raised
% translations sit clear of the line above
\newenvironment{ealglossed}{\par\linespread{2.1}\selectfont\ignorespaces}{\par}

% a whole translated sentence above its English counterpart. block level,
% so it does not belong inside a tabular cell
\newcommand{\ealpara}[2]{%
  \par\addvspace{0.5\baselineskip}%
  {\color{ealtwocolour}\ealtextblock{#1}}%
  \penalty200
  {\color{ealenglish}#2\par}%
  \addvspace{0.5\baselineskip}%
}

% ================================================================
% Dictionary layout
% ================================================================
\usepackage{multicol}

\newlength{\ealdictgap}
\setlength{\ealdictgap}{1.2mm}

% english word / english meaning / translated word / translated meaning
\newcommand{\dictentrytr}[4]{%
  \par\addvspace{\ealdictgap}%
  \noindent
  {\bfseries\ealkey{#1}}\quad{\small #2}\par
  \nopagebreak
  {\color{ealtwocolour}\ealtext{{\bfseries #3}~\textemdash{} {\small #4}}\par}%
  \par
}

\pagestyle{empty}
\setlength{\parindent}{0pt}

\begin{document}
{\large\bfseries Chinese (Traditional) Mathematical Dictionary}\par\medskip

\dictentrytr{coefficient}{the number in front of a letter, showing how many of it there are}{係數}{字母前面的數，表示有多少個該字母}
\dictentrytr{decimal}{a number written with a point, where the digits after the point are parts of a whole}{小數}{用小數點寫成的數，其中小數點後的數字是一個整體的若干部分}
\dictentrytr{denominator}{the number below the line in a fraction, showing how many equal parts the whole is split into}{分母}{分數線下方的數，表示整體被分成多少個等份}
\dictentrytr{difference}{the answer when one number is subtracted from another}{差}{一個數減去另一個數時所得的答案}
\dictentrytr{digit}{a single figure from 0 to 9 that a number is written with}{數字}{書寫一個數時所用的0至9其中一個單一數字}
\dictentrytr{equivalent}{having the same value, even though it is written differently}{等值}{即使寫法不同，仍有相同的值}
\dictentrytr{estimate}{a sensible answer worked out quickly from rounded numbers}{估算}{由捨入後的數快速算出的合理答案}
\dictentrytr{expand}{to multiply out brackets so the expression has no brackets left}{展開}{把括號乘開，使表達式不再有括號}
\dictentrytr{expression}{a group of numbers and letters joined by operations, with no equals sign}{表達式}{由運算符號連接的一組數字和字母，沒有等號}
\dictentrytr{factor}{a whole number that divides exactly into another number}{因數}{能整除另一個數的整數}
\dictentrytr{factorise}{to write an expression as a product, usually by putting brackets back in}{因式分解}{把表達式寫成乘積，通常是把括號放回去}
\dictentrytr{formula}{a rule written with letters, showing how to work one quantity out from others}{公式}{用字母寫成的規則，顯示如何從其他量算出某一個量}
\dictentrytr{fraction}{a number written as one whole split into equal parts, with a number of those parts taken}{分數}{寫成把一個整體分成等份，並取其中若干份的數}
\dictentrytr{gradient}{how steep a line is, found by how far it rises for each step across}{斜率}{一條線有多陡，由每橫向移動一步時上升多少求出}
\dictentrytr{highest common factor}{the largest whole number that divides exactly into two or more numbers}{最大公因數}{能整除兩個或以上數目的最大整數}
\dictentrytr{identity}{a statement that is true for every value of the letters in it}{恆等式}{對其中字母的每一個值都成立的陳述}
\dictentrytr{index}{the small raised number showing what power something is raised to}{指數}{顯示某數被提升到甚麼次方的小型上標數字}
\dictentrytr{inequality}{a statement that one quantity is larger or smaller than another}{不等式}{表示一個量比另一個量大或小的陳述}
\dictentrytr{integer}{a whole number, which may be positive, negative or zero}{整數}{可以是正數、負數或零的整數}
\dictentrytr{lowest common multiple}{the smallest number that two or more numbers all divide into exactly}{最小公倍數}{兩個或以上的數都能整除的最小數}
\dictentrytr{multiple}{a number you get by multiplying another number by a whole number}{倍數}{把另一個數乘以一個整數所得的數}
\dictentrytr{multiply}{to add a number to itself a given number of times}{乘}{把一個數自身相加指定次數}
\dictentrytr{numerator}{the number above the line in a fraction, showing how many equal parts are taken}{分子}{分數線上方，表示取了多少個等份的數}
\dictentrytr{percentage}{a number of parts out of a hundred}{百分數}{每一百份中的若干份}
\dictentrytr{place value}{the value a digit has because of where it sits in a number}{位值}{一個數字因它在數中的位置而具有的值}
\dictentrytr{power}{how many times a number is multiplied by itself}{冪}{一個數自乘若干次的次數}
\dictentrytr{prime}{a whole number greater than 1 that divides only by itself and 1}{質數}{大於1，且只能被自身和1整除的整數}
\dictentrytr{product}{the answer when two or more numbers are multiplied}{積}{兩個或以上的數相乘時所得的答案}
\dictentrytr{proportion}{a relationship where two amounts change by the same factor}{比例}{兩個量以相同因數改變的關係}
\dictentrytr{quotient}{the answer when one number is divided by another}{商}{一個數除以另一個數時所得的答案}
\dictentrytr{ratio}{a way of comparing two or more amounts, showing how much of one there is for each of the other}{比}{比較兩個或以上數量的方法，顯示每一個量相對於其他量有多少}
\dictentrytr{reciprocal}{the number you multiply by to get 1, found by turning a fraction upside down}{倒數}{與它相乘會得到1的數，可把分數上下倒轉求得}
\dictentrytr{recur}{to repeat the same digits forever after the decimal point}{循環}{在小數點後永遠重複相同的數字}
\dictentrytr{root}{a number that gives another number when multiplied by itself a given number of times}{根}{將自身相乘指定次數後會得到另一個數的數}
\dictentrytr{round}{to replace a number with a nearby simpler one}{捨入}{用一個附近較簡單的數取代一個數}
\dictentrytr{scale factor}{the number every length is multiplied by when a shape is enlarged}{比例因子}{圖形放大時，每一條長度所乘以的數}
\dictentrytr{simplify}{to write something in its smallest or plainest form without changing its value}{簡化}{在不改變其值的情況下，把某物寫成最小或最簡單的形式}
\dictentrytr{sum}{the answer when numbers are added together}{和}{數目相加時所得的答案}
\dictentrytr{surd}{a root that cannot be written exactly as a fraction, so it is left in root form}{不盡根}{不能準確寫成分數的根，因此保留為根式}
\dictentrytr{unitary}{a method that first finds the value of one part or one item}{歸一法}{先求出一份或一件物品的數值的方法}

\end{document}
